\documentclass[12pt,a4paper]{article}

\usepackage[utf8]{inputenc}
\usepackage[T1]{fontenc}
\usepackage[french]{babel}
\usepackage{lmodern}
\usepackage{geometry}
\usepackage{enumitem}
\usepackage{hyperref}
\usepackage{microtype}

\geometry{margin=2.5cm}
\setlength{\parskip}{0.5em}
\setlength{\parindent}{0pt}
\hypersetup{
    colorlinks=true,
    urlcolor=blue,
    linkcolor=black
}

\begin{document}

\section*{Analyse critique}

\subsection*{Limites et biais}

\begin{itemize}[leftmargin=*, itemsep=0.5em]
    \item \textbf{Biais de survie :} les cinq titres sont des gagnantes actuelles, choisies après coup. Cela gonfle l'alpha (0,76\,\% par mois) : un investisseur de 2015 n'aurait pas pu les connaître à l'avance.

    \item \textbf{Concentration :} cinq grandes valeurs technologiques, donc un risque propre au portefeuille important (erreurs de 3 à 5\,\% certains mois, section 7.3) et des conclusions non généralisables au marché.

    \item \textbf{Peu de données de test :} 33 mois. Le R² hors échantillon est bruité, et les intervalles bootstrap du CAPM et de Fama-French se chevauchent. Le gain des trois facteurs repose surtout sur le test F, calculé sur les 132 mois.

    \item \textbf{Modèle explicatif, pas prédictif :} les facteurs du mois sont connus après coup. Le modèle explique le rendement d'un mois, il ne le prévoit pas. De plus, Mkt-RF contient en partie les cinq titres.

    \item \textbf{Coefficients non constants :} le bêta SMB passe de $-0{,}53$ (2015--2019) à $-0{,}28$ (2020--2025). Trois facteurs seulement : momentum, rentabilité et investissement sont absents. Frais de transaction ignorés ; données Yahoo non officielles et sujettes à révision.
\end{itemize}

\subsection*{Pistes d'amélioration concrètes}

\begin{enumerate}[leftmargin=*, itemsep=0.5em]
    \item Élargir le portefeuille à 20 à 30 titres, en incluant des entreprises disparues ou en déclin (base CRSP), pour réduire le biais de survie et le risque spécifique.

    \item Ajouter le facteur momentum (\texttt{F-F\_Momentum\_Factor}) ou passer au modèle à 5 facteurs (\texttt{F-F\_Research\_Data\_5\_Factors\_2x3}) et comparer les R² ajustés.

    \item Estimer des bêtas variables dans le temps (fenêtre glissante de 36 mois) et calculer des erreurs-types par sous-période pour tester la stabilité.

    \item Remplacer le découpage unique par une validation à fenêtre croissante et un bootstrap par blocs, plus adapté à des séries temporelles autocorrélées.
\end{enumerate}

\section*{9. Conclusion}

Oui : pour ce portefeuille, le modèle à trois facteurs explique mieux les rendements que le CAPM seul. Son R² hors échantillon est de 0,73 contre 0,57, et le test F de SMB et HML donne $p < 0{,}0001$. L'objectif mesurable de la section 1 est atteint. Le portefeuille est exposé au marché (bêta 1,15), aux grandes capitalisations (SMB $-0{,}36$) et à la croissance (HML $-0{,}56$), et son alpha de 0,76\,\% par mois, intervalle [0,29 ; 1,24], est significatif. Ce résultat reste à nuancer : seuls 33 mois servent au test, et l'alpha est probablement gonflé par le biais de survie, donc il ne prouve pas une compétence.

\section*{10. Références et reproduction}

\subsection*{Sources de données}

\begin{itemize}[leftmargin=*, itemsep=0.5em]
    \item Kenneth R. French, \textit{Data Library} (fichier \texttt{F-F\_Research\_Data\_Factors}), Tuck School of Business, Dartmouth :
    \url{https://mba.tuck.dartmouth.edu/pages/faculty/ken.french/data_library.html}.
    Extraction : 17 novembre 2026.

    \item Yahoo Finance, cours mensuels ajustés, via la bibliothèque \texttt{yfinance} :
    \url{https://finance.yahoo.com} et
    \url{https://github.com/ranaroussi/yfinance}.
    Extraction : 17 novembre 2026.
\end{itemize}

\subsection*{Références scientifiques et logicielles}

\begin{itemize}[leftmargin=*, itemsep=0.5em]
    \item Fama, E. F., French, K. R. (1993). Common risk factors in the returns on stocks and bonds. \textit{Journal of Financial Economics}, 33(1), 3--56.

    \item Sharpe, W. F. (1964). Capital asset prices: a theory of market equilibrium under conditions of risk. \textit{The Journal of Finance}, 19(3), 425--442.

    \item Newey, W. K., West, K. D. (1987). A simple, positive semi-definite, heteroskedasticity and autocorrelation consistent covariance matrix. \textit{Econometrica}, 55(3), 703--708.

    \item Pedregosa, F. et al. (2011). Scikit-learn: machine learning in Python. \textit{Journal of Machine Learning Research}, 12, 2825--2830.

    \item Seabold, S., Perktold, J. (2010). Statsmodels: econometric and statistical modeling with Python. \textit{Proceedings of the 9th Python in Science Conference}.

    \item Support du cours : notebook « 29 exemples Python, principaux algorithmes de Machine Learning » (annexe A), algorithme n° 1 (régression linéaire).
\end{itemize}

\end{document}
